\documentclass[10pt,a4paper]{amsart}
\usepackage[margin=19mm]{geometry}
\usepackage{amsmath,amssymb,mathtools}
\usepackage[scr=rsfs,cal=euler]{mathalfa}
\usepackage{xfrac}
\usepackage[unicode,hidelinks,bookmarks=false]{hyperref}
\newcommand{\ie}{i.e.\ }
\newcommand{\cf}{cf.\ }
\newcommand{\resp}{respectively\ }
\newcommand{\Subsection}{\subsection}
\providecommand{\nobreakdash}{\nobreak\hbox{-}}
\setlength{\emergencystretch}{2em}
\multlinegap0pt
\usepackage{bm}
\usepackage{bbm}
\usepackage{dsfont}
\usepackage{xspace}
\usepackage{cancel}
\usepackage{mathtools}
\usepackage{amsthm}
\makeatletter
\@ifundefined{proposition}{\newtheorem{proposition}{Proposition}}{}
\makeatother
\title[On the rate of convergence of $\mathbb{R}^d$-ergodic average]{On the rate of convergence of $\mathbb{R}^d$-ergodic averages constructed over a strictly convex set}
\author{Ivan Podvigin}
\date{Source: arXiv:2506.16740v1 (CC BY 4.0)}
\begin{document}
\maketitle

\noindent\emph{Source and licence.} On the rate of convergence of $\mathbb{R}^d$-ergodic averages constructed over a strictly convex set. Version arXiv:2506.16740v1 (CC BY 4.0), distributed under the Creative Commons Attribution 4.0 International licence, as recorded in the arXiv licence metadata for this exact version and shown on the abstract page https://arxiv.org/abs/2506.16740v1. Full rights notice: licenses/arxiv-2506.16740-CC-BY-4.0.txt.\par\smallskip

\noindent\textit{Editorial scope.} Counted unit: the Proposition whose source label is \texttt{ProposSubHomogeneous} in the article (source0.tex lines 245--247, source label \texttt{ProposSubHomogeneous}), delivered together with the article's own complete original proof of it (source0.tex lines 249--295, one \texttt{proof} environment of 47 lines). No mathematics is written, repaired or re-derived: statement and proof are the article's own text line for line. Required context retained uncounted: none. Every definition, display and label of those lines is retained unchanged. Omitted from this package and not redistributed: 1--244, 248--248, 296--579. Nothing inside a retained excerpt is omitted or altered: every retained body is delivered exactly as the article's lines are written, so no omission receipt of any kind accompanies this package. Citations: the retained excerpts carry no citation command at all, so no bibliography entry of the article is transcribed and no reference is redistributed. Reference adaptation: none was needed and none was made; every reference inside the delivered unit resolves inside it, so no unresolved reference or citation is produced. Printed slips kept literal: no printed word, symbol, hypothesis or number of the counted statement or of its proof was altered. Layout and class: the article's own class is \texttt{article} with packages including amsthm, amsfonts, amsmath, amssymb, bm, eucal, inputenc, babel, color, graphicx; the wrapper is the lane's pinned \texttt{amsart} editorial wrapper at 10pt on A4, which restates the theorem-like environments the retained text needs and adds the article's own macro definitions that the retained text needs (none), with extra wrapper packages (bm, bbm, dsfont, xspace, cancel, mathtools). The delivered numbering is the unit's own. Assets: no retained span contains a figure environment, a graphics-inclusion command, a PDF-page inclusion, a table or a TikZ picture; no image file is copied into this package, the package declares no asset dependency and no image file is redistributed with it. Licence: version arXiv:2506.16740v1 of this article is distributed under the Creative Commons Attribution 4.0 International licence, as recorded in the arXiv licence metadata for this exact version and shown on the abstract page https://arxiv.org/abs/2506.16740v1. A full rights notice is delivered at licenses/arxiv-2506.16740-CC-BY-4.0.txt. Characters: every retained excerpt is pure ASCII, so the delivered engine, XeLaTeX with its default Latin Modern text font, reports no missing character.
\begin{proposition}\label{ProposSubHomogeneous}
Let $X\subset\mathbb{R}^d$ be a cone set. If $S_X$ is a compact set, then all homogeneous functions of the same degree defined on $X$, continuous and positive on the unit sphere $S_X$, are equivalent. Conversely, if all homogeneous functions of the same degree defined on $X$, continuous and positive on the unit sphere $S_X$, are equivalent, then $S_X$ is a compact set.
\end{proposition}

\begin{proof}
Let $\varphi$ be a homogeneous function of degree~$\theta,$ continuous and positive on the unit sphere $S_X$. Show that there exist constants~${E,F>0}$ such that for all ${{\bf t}\in X}$
$$
E|{\bf t}|^{\theta}\leq\varphi({\bf t})\leq F|{\bf t}|^{\theta}.
$$
Indeed, since
$$
\varphi({\bf t})=|{\bf t}|^\theta\varphi\left(\frac{t_1}{|t|},\frac{t_2}{|t|},...,\frac{t_d}{|t|}\right),
$$
then
$$
E=\inf\limits_{{\bf s}\in S_X}\varphi({\bf s}),\ \ F=\sup\limits_{{\bf s}\in S_X}\varphi({\bf s}).
$$
Since $\varphi$ is continuous and positive on the unit sphere~$S_X$ and $S_X$ is a compact set, then ${0<E\leq F<\infty}.$

Let us now prove the converse statement. Let ${\bf s}_n$ be a sequence in $S_X.$ There is a subsequence ${\bf s}_{n_k}$ and a vector ${\bf s}\in\mathbb{R}^d$ such that ${|{\bf s}_{n_k}-{\bf s}|\to0}$ as ${k\to\infty}.$ Suppose that ${{\bf s}\not\in S_X},$ and consequently ${r{\bf s}\not\in X}$ for all $r>0.$ 
Without loss of generality, we will assume that ${s_1\neq0}.$
Consider the homogeneous function
$$
\varphi_{\bf s}({\bf t})=\left(\sum_{j=1}^d\bigl||t_1|s_j-t_j|s_1|\bigr|\right)^\theta,\ \ {\bf t}\in X
$$
of nonzero degree $\theta.$ It is well-defined on $X,$ positive and continuous since 
$$
\sum_{j=1}^d\bigl||t_1|s_j-t_j|s_1|\bigr|=0\ \ \ \Leftrightarrow\ \ \ {\bf t}=r{\bf s}
$$
for some $r>0.$
From the equivalence condition we have 
$$
E|{\bf t}|^{\theta}\leq\varphi_{\bf s}({\bf t})\leq F|{\bf t}|^{\theta}
$$
for all ${\bf t}\in X$ and some $E,F>0.$ For ${\bf t}={\bf s}_{n_k}$ we get
$$
E\leq\varphi_{\bf s}({\bf s}_{n_k})\leq F.
$$
Passing to the limit as $k\to\infty$, we get a contradiction, since
$$
\lim_{k\to\infty}\varphi_{\bf s}({\bf s}_{n_k})=
\left\{\begin{array}{cl}
0& \text{if}\ \theta>0\\
+\infty& \text{if}\ \theta<0.
\end{array}\right.
$$
For the case ${\theta=0}$ to obtain a contradiction, it is sufficient to take the function
$$
\varphi_{\bf s}(\bf t)=\left|{\bf s}-\frac{\bf t}{|\bf t|}\right|.
$$
\end{proof}

\end{document}
